\documentclass{article}
\usepackage[utf8]{inputenc}
\usepackage[spanish]{babel}
\usepackage{amsmath}
\usepackage{graphicx}
\usepackage{booktabs}
\usepackage{float}
\usepackage{hyperref}
\usepackage{geometry}
\geometry{a4paper, margin=1in}

\begin{document}

\section{Fase de Baseline: Benchmark Inicial y Evaluación de Modelos Aislados}

\subsection{Propósito y Transición a la Experimentación Empírica}

Una vez establecidos los casos de estudio, la infraestructura de ingestión de datos y la arquitectura base de MiroFish (descritos en la sección anterior), resultó metodológicamente imperativo establecer un \textit{benchmark} riguroso utilizando modelos de lenguaje en aislamiento (paradigma \textit{Single-Agent} / \textit{Single-Model}). 

Esta fase tuvo un propósito dual: en primer lugar, cuantificar el rendimiento predictivo puro (\textit{Forecasting}) del sistema frente a los distintos escenarios (económicos, electorales y deportivos); y en segundo lugar, extraer métricas de referencia sobre el comportamiento intrínseco, los sesgos y las deficiencias de cada arquitectura LLM individual antes de introducir la complejidad combinatoria de un entorno Multi-Agente.

\subsection{Evaluación del Escalamiento Topológico (Rondas y Densidad)}

El primer parámetro de investigación consistió en medir la saturación y sensibilidad de la simulación alterando la profundidad del debate (volumen de rondas) y la densidad de conexión del grafo relacional. Estas pruebas se ejecutaron transversalmente a través de los casos de estudio, revelando que el escalamiento topológico produce efectos diametralmente opuestos dependiendo de la naturaleza determinista del entorno.

En entornos algorítmicamente cerrados y puramente cuantitativos (como el Índice de Precios al Consumidor), comprobamos una correlación directa entre el volumen de debate y la precisión inferencial. Escalar la simulación a 80 rondas con una densidad moderada (condición R80-D2) redujo el Error Absoluto Medio (MAE) a un óptimo de 1.84\%.

Sin embargo, al aplicar esta misma matriz de escalamiento a escenarios abiertos y polarizados (como el balotaje presidencial en Bolivia), observamos un ciclo destructivo. El modelo retroalimentó de forma recursiva sus propias asunciones, formando una burbuja epistemológica que lo alejó por completo del \textit{Ground Truth}.

\begin{table}[H]
    \centering
    \begin{tabular}{@{}lccccc@{}}
        \toprule
        \textbf{Entorno Evaluado} & \textbf{Variante Topológica} & \textbf{Rondas (R)} & \textbf{Densidad (D)} & \textbf{MAE Resultante} \\ \midrule
        Cuantitativo (IPC) & R10-D2 & 10 & 2 & 3.12\% \\
        Cuantitativo (IPC) & R40-D1 & 40 & 1 & 2.45\% \\
        Cuantitativo (IPC) & R40-D2 & 40 & 2 & 2.31\% \\
        \textbf{Cuantitativo (IPC)} & \textbf{R80-D2} & \textbf{80} & \textbf{2} & \textbf{1.84\% (Óptimo)} \\ \midrule
        Abierto (Electoral) & R10-D2 & 10 & 2 & 7.60 (Sesgo Inicial) \\
        \textbf{Abierto (Electoral)} & \textbf{R40-D2} & \textbf{40} & \textbf{2} & \textbf{10.0 (Herd Behavior)} \\ \bottomrule
    \end{tabular}
    \caption{Matriz de escalamiento topológico utilizando Llama-3.3-70B como ancla. El aumento de rondas mejora la precisión en entornos puramente matemáticos, pero desencadena \textit{Herd Behavior} (Cámaras de Eco) en entornos sociales.}
    \label{tab:matriz_topologica}
\end{table}

\subsection{Análisis Cruzado de Arquitecturas y Manifestación de Sesgos}

Habiendo identificado una configuración topológica de control (R40-D2), procedimos a cruzar el rendimiento predictivo y conductual de las distintas arquitecturas (Llama-3.3-70B, Gemma-3-27B y Qwen3-8B) para evaluar la varianza estadística y la propensión a sesgos de pre-entrenamiento.

Los resultados demostraron que las patologías cognitivas varían drásticamente según el modelo utilizado para orquestar la simulación:

\begin{table}[H]
    \centering
    \begin{tabular}{@{}lllc@{}}
        \toprule
        \textbf{Modelo Evaluado} & \textbf{Caso de Estudio} & \textbf{Condición de Estrés} & \textbf{Resultado Observado} \\ \midrule
        \multicolumn{4}{c}{\textit{Varianza en Precisión Predictiva (Entorno Cuantitativo)}} \\ \midrule
        Llama-3.3-70B & IPC Argentina & Réplicas de Control & MAE Medio: 1.85\% ($\sigma$: 0.05) \\
        Gemma-3-27B & IPC Argentina & Cruce Arquitectónico & MAE: 2.67\% \\
        Qwen3-8B & IPC Argentina & Cruce Arquitectónico & MAE: 3.45\% \\ \midrule
        \multicolumn{4}{c}{\textit{Vulnerabilidades en Escenarios Sociales Abiertos}} \\ \midrule
        \textbf{Llama-3.3-70B} & \textbf{Electoral (Bolivia)} & \textbf{Aislamiento Extendido} & \textbf{Burbuja Epistemológica} \\
        \textbf{Gemma-3-27B} & \textbf{Deportivo (Fútbol)} & \textbf{Inyección de Evidencia} & \textbf{Terquedad (Ignora Datos)} \\
        \textbf{Qwen3-8B} & \textbf{Multi-escenario} & \textbf{Aislamiento Inicial} & \textbf{Consenso Neutro (Pasividad)} \\ \bottomrule
    \end{tabular}
    \caption{Análisis cruzado de arquitecturas y vulnerabilidades observadas. En escenarios abiertos, Gemma demuestra impermeabilidad factual debido a su pre-entrenamiento, mientras Llama tiende a la confirmación recursiva.}
    \label{tab:evaluacion_arquitecturas}
\end{table}

Como se aprecia en la Tabla \ref{tab:evaluacion_arquitecturas}, en el escenario deportivo Gemma-3 demostró una impermeabilidad absoluta frente a la información factual (estadísticas provistas contrarias a su sesgo). El peso de su pre-entrenamiento anuló su capacidad de asimilar evidencia. Por otro lado, Qwen3 exhibió una excesiva pasividad en aislamiento, pero gran flexibilidad para absorber inyecciones de datos sorpresivos, aunque incapaz de propagar un consenso equilibrado por sí solo.

\subsection{Permeabilidad y Optimización de la Ingesta de Datos}

Otro vector fundamental del \textit{benchmark} fue evaluar los cuellos de botella en la asimilación de información. En simulaciones controladas, implementamos una optimización de deduplicación semántica previa a la inserción en el almacenamiento (Neo4j). Reducir el ruido topológico y fusionar los nodos semánticamente idénticos concentró la atención de los agentes, mejorando la precisión base de Llama-3.3 a un récord histórico de MAE 1.475\%.

No obstante, al intentar dotar al sistema de recolección de contexto 100\% autónomo (mediante \textit{Deep Search} con cortes temporales estrictos para evitar fugas de información), la precisión se degradó. Esto confirmó que, incluso en los modelos de frontera, la ausencia de anclas documentales estructuradas dispersa la atención del agente, incapacitándolo para emitir proyecciones consistentes en un entorno multi-variable.

\subsection{Sensibilidad Temporal y Robustez ante Ruido (Inyecciones Exógenas)}

Además de las pruebas topológicas y de permeabilidad pura, el \textit{benchmark} evaluó la sensibilidad del sistema frente a variaciones en la ventana temporal de los datos semilla (\textit{Seed Data}) y su respuesta ante la inyección controlada de ruido.

\subsubsection{Curva de Degradación Temporal (Ventanas de Seed Data)}
Para medir el decaimiento de la certeza predictiva, evaluamos el rendimiento en distintos horizontes temporales ($\Delta 1$ a $\Delta 4$) utilizando un \textit{cutoff} estricto en los datos semilla. El análisis reveló un patrón de degradación no lineal:
\begin{itemize}
    \item \textbf{Corto Plazo ($\Delta 1$):} El sistema exhibió una precisión excepcional (error de 0.1 puntos porcentuales) al identificar los anclajes nominales inmediatos presentes en el corpus semilla.
    \item \textbf{Medio-Corto Plazo ($\Delta 2$):} Se registró el mayor fallo analítico. Al carecer de estímulos sobre \textit{shocks} transitorios futuros, los modelos aislaron su razonamiento en una extrapolación lineal de las tendencias iniciales, subestimando severamente los rebotes macroeconómicos.
    \item \textbf{Medio Plazo ($\Delta 3$):} Irónicamente, la precisión se recuperó al evaluar tendencias macro-direccionales de fondo, probando que el motor de inferencia es robusto para la tendencia dominante pero ciego ante fluctuaciones transitorias si la ventana de inyección no se actualiza (variación de período).
\end{itemize}

\subsubsection{Inyección de Ruido y Momento de Entrada}
Adicionalmente, experimentamos alterando el momento de inyección de datos y saturando el entorno con ruido deliberado (ej. \texttt{input\_04\_noise\_dolar.txt}). Bajo condiciones de alta profundidad (R80-D2), la inyección de ruido financiero desencadenó un mecanismo inesperado: el modelo asimiló la especulación social (\textit{Herd Behavior}) de los datos ruidosos, lo que irónicamente lo llevó a predecir con extrema precisión un anclaje recesivo (logrando un récord histórico de MAE 0.9875\% en esa variante aislada). 

Sin embargo, cruzando estos datos con el análisis cualitativo (Tabla \ref{tab:evaluacion_arquitecturas}), confirmamos que el \textit{momento} de la inyección es crítico. Datos (incluso factualmente correctos) inyectados a mitad de debate suelen ser ignorados por arquitecturas con fuertes sesgos de pre-entrenamiento (como Gemma), probando que la asimilación depende del balance entre el "ruido" externo y el tiempo de arraigo de las creencias en la simulación.

\subsection{Análisis Intrínseco: Diversidad Léxica y Estabilidad Temporal}

Conscientes de las profundas patologías conductuales expuestas en el análisis cruzado, estructuramos un marco de evaluación final enfocado exclusivamente en medir el funcionamiento intrínseco de los modelos, de manera independiente a su puntería predictiva. A través de métricas algorítmicas agnósticas (entropía categórica, \textit{Self-BLEU}, \textit{Vendi Score} y deriva euclidiana), procesamos simulaciones masivas en paralelo para establecer cuánta estructura y variedad lingüística es capaz de generar cada motor.

Los resultados forenses confirmaron asimetrías severas en el comportamiento generativo:

\begin{table}[H]
    \centering
    \begin{tabular}{@{}lccccc@{}}
        \toprule
        \textbf{Modelo Evaluado} & \textbf{Posts Originales} & \textbf{Comentarios} & \textbf{Self-BLEU ($\downarrow$)} & \textbf{Distinct-2 ($\uparrow$)} & \textbf{Deriva Temporal} \\ \midrule
        Llama-3.3-70B & 10 & 57 & 0.412 & 0.614 & 0.266 \\
        \textbf{Gemma-3-27B} & \textbf{18} & \textbf{16} & \textbf{0.264*} & \textbf{0.752*} & \textbf{N/A} \\
        Qwen3-8B & N/A & N/A & N/A & N/A & \textbf{0.225*} \\ \bottomrule
    \end{tabular}
    \caption{Métricas de evaluación intrínseca. Las celdas marcadas con asterisco (*) y negrita indican la variante arquitectónica dominante: Gemma lidera contundentemente en diversidad léxica y proactividad relacional, mientras que Qwen lidera en estabilidad ideológica (menor deriva temporal).}
    \label{tab:analisis_intrinseco}
\end{table}

\begin{enumerate}
    \item \textbf{Composición Conductual y Proactividad:} La distribución de acciones en la plataforma virtual expuso un sesgo reactivo masivo en Llama-3.3: por cada 10 publicaciones originales (\textit{posts}), generó 57 comentarios a terceros. Por el contrario, Gemma-3 mostró un perfil proactivo e independiente, emitiendo 18 publicaciones y 16 comentarios.
    \item \textbf{Entropía Léxica y Repetitividad:} Al filtrar la métrica exclusivamente hacia publicaciones originales, descubrimos que Gemma-3 posee una entropía semántica significativamente mayor, produciendo textos altamente diversos (\textit{Self-BLEU} 0.264, \textit{distinct-2} 0.752). Llama-3.3 demostró una marcada pobreza en su diversidad generativa (\textit{Self-BLEU} 0.412), tendiendo a reciclar estructuras léxicas de manera sistémica.
    \item \textbf{Deriva Temporal (\textit{Intra-Persona Drift}):} Empleando encuestas paramétricas (\textit{Checkpoints Interview}) a los agentes simulados al inicio, mitad y cierre de la simulación, cuantificamos la consistencia de las posturas. Llama-3.3 sufrió de inestabilidad crónica, exhibiendo una deriva estructural mayor (\textit{Endpoint-distance} de 0.266) que Qwen3 (0.225), probando que las "personas" simuladas por Llama tienden a olvidar o mutar sus posturas originarias si la simulación se prolonga.
\end{enumerate}

\subsection{Conclusión de la Fase de Aislamiento}

La aglomeración de estos datos empíricos estableció un límite metodológico definitivo. Ningún LLM de frontera en solitario —independientemente de sus hiperparámetros o su volumen de parámetros— posee la completitud conductual necesaria para emular un foro social realista de forma ininterrumpida. La repetitividad léxica (Llama), la formación espontánea de cámaras de eco y la terquedad anclada en el pre-entrenamiento (Gemma) siempre terminan por degradar la simulación. 

Este diagnóstico exhaustivo fue el catalizador que nos forzó a descartar el enfoque de modelo único, motivando el diseño empírico de la arquitectura Multi-Agente heterogénea detallada en la sección subsecuente, bajo la premisa de que la fricción entre la proactividad de Gemma, la reactividad de Llama y la permeabilidad de Qwen actuaría como el único mecanismo de estabilización algorítmica viable.

\end{document}

